\documentclass[12pt]{article}

% --- Packages ---
\usepackage[utf8]{inputenc}
\usepackage[T1]{fontenc}
\usepackage[french]{babel}
\usepackage{amsmath, amssymb}
\usepackage{geometry}
% \usepackage[draft]{graphicx}
\usepackage{tikz}
\usepackage[most]{tcolorbox}
\usepackage{multicol}
\usepackage{lmodern}
\usepackage{pgfplots}
\pgfplotsset{compat=1.18}
\usepackage{mdframed}
\usepackage{array}
\usepackage[thinlines]{easytable}
\usepackage{tkz-tab}
\usepackage{bm}
\usepackage{xcolor}
\usepackage{sectsty}
\usepackage[utf8]{inputenc}
\usepackage[T1]{fontenc}
\usepackage[inline]{enumitem}
\usepackage[dvipsnames]{xcolor}
\usepackage{multido}
% --- Page ---
\geometry{margin=2cm}

% Footer
\usepackage{fancyhdr}
\pagestyle{fancy}
\cfoot{\thepage}
\rfoot{\textit{G.M.}}
\renewcommand{\headrulewidth}{0pt}
\renewcommand{\footrulewidth}{0pt}

% Pour les normes : 
\newcommand*{\norme}[1]{\left\lVert\overrightarrow{#1}\right\rVert}

\usetikzlibrary{shapes.misc}

\tikzset{cross/.style={cross out, draw=black, minimum size=2*(#1-\pgflinewidth), inner sep=0pt, outer sep=0pt},
%default radius will be 1pt. 
cross/.default={4pt}}

% mettre dans le préambule, après \usepackage{xcolor}
\usepackage{sectsty} % gère facilement la couleur des titres
\setlength{\parindent}{0pt}

\sectionfont{\color{red}\Large\bfseries}     % pour \section*
\subsectionfont{\color{red}\large\bfseries\hspace{1cm}}

\begin{document}

% ============================
%   Probabilités conditionnelles
% ============================

\newpage

\begin{center}
    {\LARGE \textbf{Loi de Bernoulli}}\\[4mm]
    {\large 1\up{ère} STMG}
\end{center}

\vspace{10pt}
\hrule
\vspace{0pt}

% ============================
%   PARTIE 1 : Calculs à l'aide d'un tableau croisé
% ============================

\section*{\textcolor{red}{I/ Expérience aléatoire à 2 épreuves}}

\vspace{5pt}

\begin{mdframed}[
    linewidth=1pt,
    linecolor=ForestGreen,
    innertopmargin=10pt,
    innerbottommargin=10pt
]
{\color{ForestGreen}\textbf{\underline{Définition :}}}  

Deux épreuves sont dites \textbf{indépendantes} si le résultat de l’une n’a aucune influence sur le
résultat de l’autre.

\end{mdframed}

{\color{red}\bf{\underline{Exemple et méthode :}} Calculer une probabilité associée à une expérience aléatoire 
à deux épreuves.}

\vspace{5pt}
On dispose dans une urne de 4 boules : 2 rouges, 1 jaune et 1 verte. Dans une seconde urne, on dispose de 
3 boules noires et 1 boule blanche.

\vspace{5pt}
Alban tire successivement 1 boule dans la première urne puis une boule dans la deuxième.

\vspace{10pt}
Ces deux épreuves, sont dites \_\_\_\_\_\_\_\_\_\_\_\_.

\vspace{10pt}
\hspace{1cm} \textbf{1) Représenter la situation à l'aide d'un arbre pondéré.}

\hspace{1cm} \textbf{2) Calculer la probabilité $P_1$ d'obtenir une boule rouge et une boule noire.}

\hspace{1cm} \textbf{3) Calculer la probabilité $P_2$ de ne \underline{pas} obtenir une boule verte et 
d'obtenir une boule blanche.}

\begin{minipage}[t]{.55\textwidth}%
    \vspace{0pt}
    \vspace{0.5cm}
    {\color{black!50}\dotfill}
    
    \vspace{0.5cm}
    {\color{black!50}\dotfill}
    
    \vspace{0.5cm}
    {\color{black!50}\dotfill}
    
    \vspace{0.5cm}
    {\color{black!50}\dotfill}
    
    \vspace{0.5cm}
    {\color{black!50}\dotfill}

    \vspace{0.5cm}
    {\color{black!50}\dotfill}
    
    \vspace{0.5cm}
    {\color{black!50}\dotfill}
    
    \vspace{0.5cm}
    {\color{black!50}\dotfill}
    
    \vspace{0.5cm}
    {\color{black!50}\dotfill}
    
    \vspace{0.5cm}
    {\color{black!50}\dotfill}
    
    \vspace{0.5cm}
    {\color{black!50}\dotfill}
    
    \vspace{0.5cm}
    {\color{black!50}\dotfill}
    
    

\end{minipage}%
\hfill
\begin{minipage}[t]{.45\textwidth}%
    \vspace{0pt}
    
\end{minipage}%

\begin{mdframed}[
    linewidth=1pt,
    roundcorner=5pt,
    linecolor=RoyalBlue,
    innertopmargin=10pt,
    innerbottommargin=10pt
]
{\color{RoyalBlue}
\textbf{\underline{Exercice d'application n°1:} Exploiter un arbre de probabilités}}

\vspace{5pt}
On s'intéresse aux personnes ayant les yeux bleus. On note $B$ l'évènement « Avoir l'allèle des yeux 
bleus ». On choisit au hasard deux personnes dans la population et on regarde si elles possèdent cet allèle. 
On suppose que 40\% de la population possède l'allèle des yeux bleus.

\vspace{5pt}
\hspace{1cm}1) La situation correspond-elle à deux épreuves indépendantes ?

\vspace{5pt}
\hspace{1cm}2) Dessiner un arbre de probabilités qui représente la situation.

\vspace{5pt}
\hspace{1cm}3) Déterminer la probabilité que les deux personnes aient l'allèle des yeux bleus.

\begin{minipage}[t]{.55\textwidth}%
    \vspace{0pt}
    \vspace{0.5cm}
    {\color{black!40}\dotfill}

    \vspace{0.5cm}
    {\color{black!40}\dotfill}

    \vspace{0.5cm}
    {\color{black!40}\dotfill}

    \vspace{0.5cm}
    {\color{black!40}\dotfill}

    \vspace{0.5cm}
    {\color{black!40}\dotfill}

    \vspace{0.5cm}
    {\color{black!40}\dotfill}

    \vspace{0.5cm}
    {\color{black!40}\dotfill}
    

    \vspace{20pt}
\end{minipage}%
\hfill
\begin{minipage}[t]{.45\textwidth}%
    \vspace{0pt}
    
\end{minipage}%

\end{mdframed}

\section*{\textcolor{red}{II/ Épreuve de Bernoulli }}

\subsection*{\hspace{1cm} 1) Définition }

\begin{mdframed}[
    linewidth=1pt,
    roundcorner=5pt,
    linecolor=ForestGreen,
    innertopmargin=10pt,
    innerbottommargin=10pt
]
{\color{ForestGreen}
\textbf{\underline{Définition :}}}

\vspace{5pt}
Une \textbf{épreuve de Bernoulli} est une expérience aléatoire à \underline{deux issues} que l'on peut 
nommer « succès » 
ou « échec ».

\end{mdframed}

\underline{Exemple :}

\begin{itemize}
    \item Le jeu du pile ou face : On considère par exemple comme succès "obtenir pile" et
    comme échec "obtenir face".
    \item  On lance un dé et on considère par exemple comme succès "obtenir un 1" et
    comme échec "ne pas obtenir un 1".

    La loi de Bernoulli associée à cette expérience est :

    \vspace{-5pt}
    \begin{center}
        \begin{tabular}{|>{\centering\arraybackslash}m{3cm}|>{\centering\arraybackslash}m{3cm}|
            >{\centering\arraybackslash}m{3cm}|}
            \hline
            \parbox[c][1cm][c]{3cm}{\centering \textbf{Issues}} & 
            \parbox[c][1cm][c]{3cm}{\centering \textbf{Succès ( S )}} & 
            \parbox[c][1cm][c]{3cm}{\centering \textbf{Échec ( $\bm{\overline{S}}$ )}} \\
            \hline
            \parbox[c][1.5cm][c]{3cm}{\centering \textbf{Probabilités}} & 
            \parbox[c][1.5cm][c]{3cm}{\centering } 
            & \\
            \hline
        \end{tabular}
    \end{center}
\end{itemize}

\begin{mdframed}[
    linewidth=1pt,
    roundcorner=5pt,
    linecolor=Red,
    innertopmargin=10pt,
    innerbottommargin=10pt
]
{\color{red}
\textbf{\underline{Propriété :}}

\vspace{5pt}
Une \textbf{loi de Bernoulli} est la loi de probabilité d'une épreuve de Bernoulli qui suit le schéma suivant :

\begin{itemize}
    \item la probabilité d'obtenir 1 (ou succès) est égale à $p$
    \item la probabilité d'obtenir 0 (ou échec) est égale à $1-p$
\end{itemize}

$p$ est appelé le \textbf{paramètre} de la loi de Bernoulli.
}
\end{mdframed}

On peut résumer la loi de Bernoulli de paramètre $p$ dans ce tableau : 

\begin{center}
    \begin{tabular}{|>{\centering\arraybackslash}m{3cm}|>{\centering\arraybackslash}m{3cm}|
        >{\centering\arraybackslash}m{3cm}|}
        \hline
        \parbox[c][1cm][c]{3cm}{\centering \textbf{$\bm{x_i}$}} & 
        \parbox[c][1cm][c]{3cm}{\centering \textbf{1 (ou S)}} & 
        \parbox[c][1cm][c]{3cm}{\centering \textbf{0 (ou $\bm{\overline{S}}$)}} \\
        \hline
        \parbox[c][1.5cm][c]{3cm}{\centering \textbf{$\bm{P(X = x_i)}$}} &
        \parbox[c][1.5cm][c]{3cm}{\centering }
        & \\
        \hline
    \end{tabular}
\end{center}

\underline{En reprenant les exemples plus haut, déterminer $p$ : }

\begin{itemize}
    \item Le jeu du pile ou face : On considère par exemple comme succès "obtenir pile" et
    comme échec "obtenir face".

    \vspace{-10pt}
    \begin{center}
        $p =$
    \end{center}

    \vspace{5pt}
    \item  On lance un dé et on considère par exemple comme succès "obtenir un 1" et
    comme échec "ne pas obtenir un 1".

    \vspace{-10pt}
    \begin{center}
        $p =$
    \end{center}
\end{itemize}

\begin{mdframed}[
    linewidth=1pt,
    roundcorner=5pt,
    linecolor=RoyalBlue,
    innertopmargin=10pt,
    innerbottommargin=10pt
]
{\color{RoyalBlue}
\textbf{\underline{Exercice d'application n°2:} Savoir reconnaître une situation modélisée par une loi 
de Bernoulli.}}

\vspace{5pt}
Après la correction d’un contrôle, le professeur compte que 24 élèves ont obtenu une note
supérieure ou égale à 10, et 6 ne l’ont pas obtenue.
Le professeur choisit une copie au hasard.

\hspace{1cm}1) Cette situation peut-elle être modélisée par une loi de Bernoulli ?

\hspace{1cm}2) Si oui, donner son paramètre.

\vspace{0.5cm}
{\color{black!40}\dotfill}

\vspace{0.5cm}
{\color{black!40}\dotfill}

\vspace{0.5cm}
{\color{black!40}\dotfill}

\vspace{0.5cm}
{\color{black!40}\dotfill}

\vspace{0.5cm}
{\color{black!40}\dotfill}


\vspace{20pt}

\end{mdframed}

\vspace{10pt}
\subsection*{\hspace{1cm} 2 Répétition d'épreuves de Bernoulli) }

\vspace{5pt}
\begin{mdframed}[
    linewidth=1pt,
    roundcorner=5pt,
    linecolor=ForestGreen,
    innertopmargin=10pt,
    innerbottommargin=10pt
]
{\color{ForestGreen}
\textbf{\underline{Définition :}}}

\vspace{5pt}
Plusieurs expériences sont identiques et indépendantes si : 
\begin{itemize}
    \item elles ont les mêmes issues,
    \item chaque issue possède la même probabilité.
\end{itemize}

\end{mdframed}

\vspace{5pt}
{\color{RedOrange}\textbf{\underline{Méthode :} Calculer une probabilité associé à une succesion d'épreuves 
de Bernoulli.}}

\vspace{5pt}
On considère l'expérience suivante :

\vspace{5pt}
Une urne contient 3 boules blanches et 2 boules rouges. On tire au hasard une boule et on la
remet dans l'urne. On considère comme succès : « obtenir une boule blanche ».

On répète l'expérience trois fois de suite.

\vspace{5pt}
\hspace{1cm}\textbf{1) A-t-on bien une situation de répétition d'épreuves de Bernoulli ? De quel paramètre ?}

\vspace{5pt}
\hspace{1cm}\textbf{2) À l’aide d’un arbre pondéré, calculer la probabilité d’obtenir deux succès.}

\vspace{0.5cm}
{\color{black!40}\dotfill}

\vspace{0.5cm}
{\color{black!40}\dotfill}

\vspace{0.5cm}
{\color{black!40}\dotfill}

\vspace{0.5cm}
{\color{black!40}\dotfill}

\vspace{0.5cm}
{\color{black!40}\dotfill}

\vspace{0.5cm}
{\color{black!40}\dotfill}

\vspace{0.5cm}
{\color{black!40}\dotfill}

\vspace{0.5cm}
{\color{black!40}\dotfill}

\vspace{0.5cm}
{\color{black!40}\dotfill}

\vspace{0.5cm}
{\color{black!40}\dotfill}

\vspace{20pt}

\newpage
\begin{mdframed}[
    linewidth=1pt,
    roundcorner=5pt,
    linecolor=RoyalBlue,
    innertopmargin=10pt,
    innerbottommargin=10pt
]
{\color{RoyalBlue}
\textbf{\underline{Exercice d'application n°3:} Calculer une probabilité associé à une succesion}

\textbf{d'épreuves 
de Bernoulli.}}

\vspace{5pt}
On lance 4 fois de suite un dé à six faces et on note à chaque fois le résultat.
À chaque lancer, on considère comme succès "obtenir un six" et comme échec "ne pas
obtenir un six".

On répète ainsi 5 fois de suite la même expérience de Bernoulli (lancer un dé) et les
expériences sont indépendantes l'une de l'autre : un lancer n'influence pas le résultat d'un
autre lancer.

\vspace{10pt}
\hspace{1cm}1) Calculer la probabilité d'obtenir 2 six.

\vspace{0.5cm}
{\color{black!45}\dotfill}

\vspace{0.5cm}
{\color{black!45}\dotfill}

\vspace{0.5cm}
{\color{black!45}\dotfill}

\vspace{0.5cm}
{\color{black!45}\dotfill}

\vspace{0.5cm}
{\color{black!45}\dotfill}

\vspace{0.5cm}
{\color{black!45}\dotfill}

\vspace{0.5cm}
{\color{black!45}\dotfill}

\vspace{0.5cm}
{\color{black!45}\dotfill}

\vspace{0.5cm}
{\color{black!45}\dotfill}

\vspace{20pt}

\end{mdframed}

\vspace{5pt}
\subsection*{\hspace{1cm} 3) Espérance }

\begin{mdframed}[
    linewidth=1pt,
    roundcorner=5pt,
    linecolor=Red,
    innertopmargin=10pt,
    innerbottommargin=-5pt
]
{\color{red}
\textbf{\underline{Propriété :}}

\vspace{5pt}
Soit $X$ une variable aléatoire qui suit une loi de Bernoulli de paramètre $p$. Alors : 
\[
E(X) = p
\]
}
\end{mdframed}






\end{document}
